\section*{XÂY DỰNG CAMPUS NETWORK\\SỬ DỤNG SDN VÀ SD-WAN\\TÓM TẮT}

\begin{itemize}
    \item[-] \textbf{Vấn đề nghiên cứu:}

    Khi trường đại học mở rộng thêm cơ sở tại nhiều địa phương, hệ thống mạng phải đồng thời bảo đảm hiệu năng, tính sẵn sàng, bảo mật, quản trị tập trung và kết nối ổn định giữa các cơ sở. Mạng truyền thống được cấu hình riêng trên từng thiết bị, thiếu cái nhìn toàn cục về trạng thái mạng và phụ thuộc một đường truyền WAN, nên mỗi thay đổi VLAN hay chính sách đều tốn công, dễ sai lệch và không tự chọn đường theo chất lượng. Đề tài thiết kế, triển khai và kiểm chứng một mô hình mạng Campus đa cơ sở kết hợp mô hình phân cấp truyền thống với SDN và SD-WAN. Báo cáo gồm 4 chương:
        \begin{itemize}
        \item[•] Chương 1 --- Mở đầu và tổng quan đề tài
        \item[•] Chương 2 --- Cơ sở lý thuyết
        \item[•] Chương 3 --- Mô hình mạng Campus đề xuất
        \item[•] Chương 4 --- Triển khai, đánh giá và kết luận
        \end{itemize}

    \item[-] \textbf{Hướng tiếp cận:}
        \begin{itemize}
        \item[•] Lý thuyết: nghiên cứu mô hình Core--Distribution--Access, VLAN, RSTP, EtherChannel, VRRP, OSPF, BGP, tường lửa trạng thái và DMZ; SDN với OpenFlow 1.3, OVS, Ryu; kiến trúc Cisco SD-WAN (Viptela) với OMP, TLOC, BFD, PKI và chính sách tập trung.
        \item[•] Thực nghiệm: xây dựng hệ thống 74 node trên EVE-NG gồm Campus chính (Site 100), ba cơ sở phụ Cần Thơ, Đà Nẵng, Nha Trang (Site 200, 300, 400), cơ sở mở rộng (Site 500), trung tâm điều khiển SD-WAN (Site 900) và hai nhà cung cấp Internet, MPLS; kiểm thử theo lớp và tự động hóa các phép đo.
        \end{itemize}

    \item[-] \textbf{Cách giải quyết vấn đề:}

    Campus chính được thiết kế ba tầng: hai Core Switch làm gateway VRRP, chạy OSPF và DHCP Relay; cặp tường lửa ASAv chạy HA bảo vệ DMZ (Web, Mail) và Server Farm (DHCP, DNS, Syslog). Sáu switch lớp 2 là OVS do ứng dụng Ryu của nhóm điều khiển: học MAC theo VLAN, cây chuyển tiếp tính tập trung kèm thăm dò liên kết chủ động, kênh quản trị VLAN 99 nhiều đường không phụ thuộc STP, tự phục hồi khi mất liên kết hoặc thiết bị, thêm VLAN và áp chính sách qua API, cùng giao diện giám sát Campus SDN Console. Các cơ sở phụ dùng mô hình Firewall-as-Core, tường lửa chi nhánh làm gateway và cấp DHCP. Các cơ sở được kết nối bằng Cisco SD-WAN với mười vEdge trên hai transport, underlay BGP, xác thực PKI, overlay IPsec và chính sách tập trung trên vSmart (SLA, AAR, data policy, control policy ưu tiên vEdge1); tường lửa chi nhánh theo dõi đường MPLS bằng IP SLA để tự chuyển sang Internet. Người dùng ra Internet tại chỗ (DIA) có dự phòng qua MPLS; vManage và NetOps Bot (Telegram) cảnh báo, chẩn đoán sự cố.

    \item[-] \textbf{Kết quả đạt được:}

    Cả 30 kịch bản kiểm thử đều đạt. Thêm VLAN qua API mất 3,8~ms, áp chính sách toàn campus 8,0--12,6~ms; sáu kịch bản cắt liên kết SDN tự phục hồi với gián đoạn không quá 3,6 giây. Mười vEdge lên fabric với toàn bộ phiên BFD hoạt động (112 đường hầm, khả dụng 99,957\%); chính sách tập trung đạt 10/10 phép thử tại cả bốn chi nhánh; qua 46 lần cắt thật, khi mất MPLS, Internet hoặc một vEdge, luồng mới tự khôi phục sau 3,2--13,4 giây; mất controller SD-WAN không làm mất gói. Dưới tải cao, OVS đạt khoảng 1,2~Gbit/s, liên VLAN khoảng 170~Mbit/s và Controller xử lý 1000 MAC mới/giây mà không ảnh hưởng luồng đang chạy. Hạn chế còn lại gồm Controller chưa có dự phòng và hiệu năng bị giới hạn bởi môi trường ảo hóa. Mô hình chứng minh SDN và SD-WAN kết hợp được với các cơ chế truyền thống trong một hệ thống thống nhất.
\end{itemize}

\newpage
